% ======================================================================
%  DepTGL 论文草稿 —— 依据 示例.tex(VLDB / ACM acmart 模板)规范化
%  只填写:摘要 v4 + 引言 8 段(均为 2026-08-28 讨论后的最新版)
%  内容来源:
%    - 论文/draft_v2_摘要.md(摘要 v4,两处挂账同步点见文内 TODO 注释)
%    - 论文/draft_v2_引言.md(两轮逐句批改后)
%  占位清单与使用说明见 README.md
% ======================================================================
% ---------- 1. 选择文档格式 ----------
% sigconf：ACM 会议标准格式（双栏，10pt 字）
% nonacm：告诉模版我们不是 ACM 的正式出版物，而是 VLDB（PVLDB 系列）
% 这一行绝对不能改！
\documentclass[sigconf, nonacm]{acmart}

%% ---------- 2. VLDB 需要的变量（投稿时可以暂时不改，最终版再填）----------
% ---- 论文专用信息 ----
% 你的论文 DOI（如果没有就保持 XX.XX/XXX.XX）
\newcommand\vldbdoi{XX.XX/XXX.XX}
% 你的论文在 PVLDB 卷中的页码（录用后会通知，投稿时留 XXX-XXX）
\newcommand\vldbpages{XXX-XXX}
\newcommand\vldbvolume{19}        % 卷号（当前 2025-2026 卷；示例里的 14 已过时，投稿时以 VLDB 官网为准）
\newcommand\vldbissue{1}          % 期号
\newcommand\vldbyear{2026}        % 年份
% ---- 这两个变量会自动取你的作者和短标题，无需改动 ----
% ---- 作品公开链接（可选） ----
% 没有公开代码/数据就留空（示例里的 URL_TO_YOUR_ARTIFACTS 会打印到页脚，已清空）
\newcommand\vldbavailabilityurl{}
% ---- 页码样式 ----
% plain：审稿版显示页码（方便审稿人引用）
% empty：最终版不显示页码
\newcommand\vldbpagestyle{plain}
% ========== 3. 文档正式开始 ==========
\begin{document}
% ---------- 3.1 标题 ----------
% TODO: DepTGL 全称待确认（待确认问题 #1）
\title{DepTGL: A NAS System for Temporal Graph Neural Networks}

%% ---------- 3.2 作者和单位 ----------
\author{Haoyu Wang}
\affiliation{%
  \institution{SDU}
  \city{ShanDong}
  \country{China}
}
\email{1778942923@qq.com}

%% ---------- 3.3 摘要 ----------
\begin{abstract}
  Temporal Graph Neural Networks (TGNNs) such as JODIE have achieved
  state-of-the-art performance on temporal interaction prediction tasks.
  However, designing TGNN architectures requires labor-intensive manual
  tuning over a large hyperparameter space, and existing neural
  architecture search (NAS) frameworks target static GNNs, whose training
  assumes i.i.d. samples and thus cannot handle the temporal data
  dependencies inherent in interaction streams. Searching over hundreds
  of candidate architectures demands parallel execution, yet naively
  parallelizing TGNN training breaks the read-after-write (RAW)
  dependencies among consecutive interactions, yielding biased
  architecture scores that cause NAS to select inferior architectures.
  % TODO(挂账同步点 1,等批改摘要时定):"among consecutive interactions" → "among interactions"
  We introduce \texttt{DepTGL}, an NAS system for TGNNs built on three
  techniques: (i) faithful execution, where parallel training respects
  the stream's RAW dependencies and a random-state protocol --- per-trial
  seed discipline, RNG-preserving state migration, and off-policy
  controller updates --- makes every backend reproduce the serial
  search's scores and selection;
  % TODO(挂账同步点 2,等批改摘要时定):"scores and selection" → "selection"
  (ii) a pipeline strategy that
  partitions the interaction stream into cost-balanced stages while
  preserving RAW semantics at stage boundaries; and (iii) asynchronous
  architecture generation that overlaps candidate generation with
  training in a persistent worker pool. Extensive experiments on X
  datasets with Y model families show that DepTGL achieves up to
  3.0$\times$ speedup over its synchronous counterpart without degrading
  search accuracy.
  % TODO: X datasets / Y model families ← 完整实验表（阶段3）
\end{abstract}
%% ---------- 3.4 生成标题区域 ----------
%% 这行命令一定要放在 \begin{abstract} 之后，它会输出标题、作者、摘要
\maketitle
%% ========== 4. VLDB 自动生成的页脚信息（绝对不要修改！） ==========
%%% do not modify the following VLDB block %%%
%%% VLDB block start %%%
\pagestyle{\vldbpagestyle}
\begingroup\small\noindent\raggedright\textbf{PVLDB Reference Format:}\\
\vldbauthors. \vldbtitle. PVLDB, \vldbvolume(\vldbissue): \vldbpages, \vldbyear.\\
\href{https://doi.org/\vldbdoi}{doi:\vldbdoi}
\endgroup
\begingroup
\renewcommand\thefootnote{}\footnote{\noindent
This work is licensed under the Creative Commons BY-NC-ND 4.0 International License. Visit \url{https://creativecommons.org/licenses/by-nc-nd/4.0/} to view a copy of this license. For any use beyond those covered by this license, obtain permission by emailing \href{mailto:info@vldb.org}{info@vldb.org}. Copyright is held by the owner/author(s). Publication rights licensed to the VLDB Endowment. \\
\raggedright Proceedings of the VLDB Endowment, Vol. \vldbvolume, No. \vldbissue\ %
ISSN 2150-8097. \\
\href{https://doi.org/\vldbdoi}{doi:\vldbdoi} \\
}\addtocounter{footnote}{-1}\endgroup
%%% VLDB block end %%%
%%% 如果设置了作品公开链接，这里会自动显示
%%% do not modify the following VLDB block %%%
%%% VLDB block start %%%
\ifdefempty{\vldbavailabilityurl}{}{
\vspace{.3cm}
\begingroup\small\noindent\raggedright\textbf{PVLDB Artifact Availability:}\\
The source code, data, and/or other artifacts have been made available at \url{\vldbavailabilityurl}.
\endgroup
}
%%% VLDB block end %%%
%% ========== 5. 正文 ==========
%% 本版只填写摘要 + 引言；第 2-7 节骨架待阶段 2-3 逐节补。
\section{Introduction}

Temporal interaction data are ubiquitous in services such as
recommendation systems, social networks, and question-answering
platforms. Each record is a (user, item, timestamp) event. Temporal
Graph Neural Networks (TGNNs) such as JODIE \cite{jodie} and TGN
\cite{tgn} achieve state-of-the-art performance for future-interaction
prediction. As each interaction arrives, they recursively update the
embeddings of the participating user and item, so a node's latest
embedding summarizes its interaction history.

However, choosing a TGNN architecture is far from trivial. A TGNN
exposes a large hyperparameter space (aggregation functions, temporal
decay functions, memory cells, time projection), and manual tuning is
labor-intensive and error-prone. Neural Architecture Search (NAS)
automates architecture design. For static GNNs, NAS frameworks such as
GraphNAS \cite{graphnas} find architectures that rival or outperform
hand-designed ones. Yet these frameworks cannot be directly applied to
temporal graphs. Static GNN training assumes i.i.d. samples, so
mini-batches can be processed in arbitrary order and in parallel. TGNN
training is different: each interaction rewrites the embeddings of its
user and item, and subsequent interactions must read the values just
written, creating read-after-write (RAW) dependencies along the
stream. TGNN training is thus inherently sequential. Since NAS must
evaluate hundreds or thousands of candidates, serial training is
prohibitively slow and parallel execution is essential. The central
challenge is to enable parallel TGNN training that respects RAW
dependencies and does not otherwise bias architecture evaluation.

Several frameworks support TGNN training. Single-machine frameworks
such as PyGT \cite{pygt}, CacheG \cite{cacheg}, and PiPAD \cite{pipad}
optimize caching, embedding reuse, and pipelined parallelism, but are
limited by the resources of a single machine. Distributed frameworks
such as ESDG \cite{esdg} and DynaHB \cite{dynahb} scale TGNN training
across multiple machines through partitioning and communication
avoidance. These systems, however, train and evaluate a fixed
architecture; they do not compare candidates. NAS, in contrast, must
evaluate hundreds or thousands of candidates, and none of these
systems provides a search strategy or evaluation mechanism for this
purpose. To the best of our knowledge, no NAS framework exists for
temporal GNNs. The natural approach, then, is to parallelize candidate
evaluation on top of existing training frameworks. As we demonstrate
next, doing so naively violates the read-after-write (RAW)
dependencies of temporal training. The resulting scores are so biased
that NAS selects the wrong architecture.

This failure is not hypothetical. On the MOOC dataset \cite{mooc}, we
ran the same NAS search three times. All three runs used the same
three-worker parallel search framework; the only variable was how each
candidate's training forms its batches: serial processing (one
interaction at a time, in stream order), conflict-free batching (each
batch contains unique nodes), or naive batching (consecutive
interactions are chunked without conflict resolution). The serial
search converges to a compact architecture with 133K parameters and a
test MRR of 0.8561. The conflict-free search selects the same
architecture, confirming that batching itself is harmless. The naive
search instead selects a 147K architecture that relies on static
embeddings; its serial re-training achieves a test MRR of only 0.6014,
a drop of 0.2547. Consider a batch of consecutive interactions
$(u_1,i_1)$, $(u_1,i_2)$, $(u_2,i_1)$. Naive batching processes the
first two against the same pre-batch embedding, so $(u_1,i_2)$ misses
the update that serial training would have applied first. The RAW
dependency along the stream is violated \cite{jodie}.
% TODO: "violated [ref]" 的引用出处按引用策略定（暂用 JODIE，冲突概念出处）
Worse, the failure is self-concealing: under naive evaluation the same
architecture scores 0.96 on validation, whereas faithful evaluation
yields 0.62, and the test MRR it reports under its own naive
evaluation, 0.9335, even exceeds the serial search's 0.8561.
Predictions computed from stale states systematically overrate
architectures that rely on static features, distort the leaderboard,
and drive NAS toward the wrong architecture.

\paragraph{Challenge I: Temporal Data Dependency.}
Interactions that touch the same node form read-after-write (RAW)
dependencies because training rewrites node embeddings as interactions
arrive. Static NAS frameworks assume i.i.d. samples and process
batches in arbitrary order; applied to TGNNs, out-of-order or
interleaved batches read stale states and produce scores that do not
reflect a candidate's true quality. Existing conflict-free batching
speeds up training but does not guarantee faithful scores --- an open
problem in NAS evaluation.

\paragraph{Challenge II: Evaluation Fidelity.}
NAS compares architectures by score, so the relative ranking of
candidates must not depend on which backend evaluates them. Parallel
backends evaluate candidates in different processes and pipeline
stages and migrate model and optimizer state across process
boundaries. Small deviations in random state and update cadence
accumulate into systematic score bias, distorting the leaderboard and
causing search to converge to an architecture favored by the
evaluation artifact rather than by the data.

\paragraph{Challenge III: Skewed Workloads.}
Temporal partitions are structurally uneven: the numbers of new users
per partition, and hence the cost of processing each partition, vary
widely across partitions (by up to 4$\times$ on MOOC), causing static
worker allocation to leave some workers idle while others become
stragglers. An evaluator must balance stages cheaply, without an
expensive profiling pass per candidate.
% 数字来源(2026-09-03, partition_stats.py):new users 每分区 max/min =
% 1.83× (20K) / 4.09× (100K) / 114× (全量);事件数严格均匀(1.00, count 策略),
% 故原 "interaction counts vary widely" 表述已删;句内 4× 锚定 100K。

To address these challenges, we present \texttt{DepTGL}, an NAS
system for temporal GNNs. DepTGL searches a JODIE-family architecture
space using a REINFORCE controller and evaluates candidates through
interchangeable execution backends (serial, data-parallel, pipeline,
asynchronous) that share a common evaluation interface. Three
techniques address these challenges while improving efficiency. First,
\textbf{faithful execution} addresses Challenges I and II. Parallel
training respects the stream's RAW dependencies, and a random-state
protocol enforces evaluation fidelity by construction. As a result,
every backend reproduces the serial search's selection. Second, a
\textbf{pipeline strategy} addresses Challenge III by partitioning the
interaction stream into stages and balancing them according to
estimated partition cost. This keeps workers saturated under skewed
workloads without violating RAW semantics at stage boundaries. Third,
an \textbf{asynchronous generation engine} maintains a persistent
worker pool in which candidate generation overlaps candidate training,
so GPUs rarely idle and the search completes in a fraction of the
serial wall-clock time.

In summary, this paper makes the following contributions:
\begin{itemize}
\item \textbf{A NAS system for temporal GNNs.} \texttt{DepTGL} is an
end-to-end framework that searches a JODIE-family architecture space
using a REINFORCE controller and evaluates candidates through
interchangeable execution backends (serial, data-parallel, pipeline,
asynchronous). Faithful execution is a design property: parallel
training respects RAW dependencies, and a random-state protocol
ensures that every backend reproduces the serial search's selection
(Challenges I and II).
\item \textbf{A pipeline strategy} that partitions the interaction
stream into cost-balanced stages and preserves RAW semantics at stage
boundaries, keeping workers saturated under skewed workloads
(Challenge III).
\item \textbf{An asynchronous generation engine} that overlaps
candidate generation with training in a persistent worker pool,
achieving a 3.0$\times$ speedup over the synchronous pipeline without
degrading search accuracy.
\end{itemize}

The remainder of this paper is organized as follows. Section 2 reviews
related work. Section 3 introduces preliminaries on temporal graphs,
JODIE-style training, and NAS. Section 4 presents an overview of
\texttt{DepTGL}. Section 5 describes faithful parallel execution,
including the pipeline strategy and asynchronous generation engine.
Section 6 reports experimental results, and Section 7 concludes.

% ========== 6. 参考文献 ==========
% TODO: 补 sample.bib（9 条:jodie / tgn / graphnas / pygt / cacheg / pipad /
%       esdg / dynahb / mooc）。未补之前文中 \cite 显示为 [?]，属正常。
% TODO: 致谢（\begin{acks}...\end{acks}）待基金信息确定后补。
\bibliographystyle{ACM-Reference-Format}
\bibliography{sample}

\end{document} 
\endinput
